\documentclass[10pt,a4paper]{amsart}
\usepackage{amsmath,amssymb,amsthm,graphicx}
\usepackage{verbatim}
\newenvironment{reference}{\comment}{\endcomment}
%\newenvironment{reference}{}{}
\newenvironment{slogan}{\comment}{\endcomment}
\newenvironment{history}{\comment}{\endcomment}

% For commutative diagrams we use Xy-pic
\usepackage[all]{xy}

% We use 2cell for 2-commutative diagrams.
\xyoption{2cell}
% Theorem environments.
%
\theoremstyle{plain}
\newtheorem{theorem}[subsection]{Theorem}
\newtheorem{proposition}[subsection]{Proposition}
\newtheorem{lemma}[subsection]{Lemma}

\theoremstyle{definition}
\newtheorem{definition}[subsection]{Definition}
\newtheorem{example}[subsection]{Example}
\newtheorem{exercise}[subsection]{Exercise}
\newtheorem{situation}[subsection]{Situation}

\theoremstyle{remark}
\newtheorem{remark}[subsection]{Remark}
\newtheorem{remarks}[subsection]{Remarks}

\numberwithin{equation}{subsection}

% Macros
%
\def\lim{\mathop{\mathrm{lim}}\nolimits}
\def\colim{\mathop{\mathrm{colim}}\nolimits}
\def\Spec{\mathop{\mathrm{Spec}}}
\def\Hom{\mathop{\mathrm{Hom}}\nolimits}
\def\Ext{\mathop{\mathrm{Ext}}\nolimits}
\def\SheafHom{\mathop{\mathcal{H}\!\mathit{om}}\nolimits}
\def\SheafExt{\mathop{\mathcal{E}\!\mathit{xt}}\nolimits}
\def\Sch{\mathit{Sch}}
\def\Mor{\mathop{\mathrm{Mor}}\nolimits}
\def\Ob{\mathop{\mathrm{Ob}}\nolimits}
\def\Sh{\mathop{\mathit{Sh}}\nolimits}
\def\NL{\mathop{N\!L}\nolimits}
\def\CH{\mathop{\mathrm{CH}}\nolimits}
\def\proetale{{pro\text{-}\acute{e}tale}}
\def\etale{{\acute{e}tale}}
\def\QCoh{\mathit{QCoh}}
\def\Ker{\mathop{\mathrm{Ker}}}
\def\Im{\mathop{\mathrm{Im}}}
\def\Coker{\mathop{\mathrm{Coker}}}
\def\Coim{\mathop{\mathrm{Coim}}}

% Boxtimes
%
\DeclareMathSymbol{\boxtimes}{\mathbin}{AMSa}{"02}

%
% Macros for moduli stacks/spaces
%
\def\QCohstack{\mathcal{QC}\!\mathit{oh}}
\def\Cohstack{\mathcal{C}\!\mathit{oh}}
\def\Spacesstack{\mathcal{S}\!\mathit{paces}}
\def\Quotfunctor{\mathrm{Quot}}
\def\Hilbfunctor{\mathrm{Hilb}}
\def\Curvesstack{\mathcal{C}\!\mathit{urves}}
\def\Polarizedstack{\mathcal{P}\!\mathit{olarized}}
\def\Complexesstack{\mathcal{C}\!\mathit{omplexes}}
% \Pic is the operator that assigns to X its picard group, usage \Pic(X)
% \Picardstack_{X/B} denotes the Picard stack of X over B
% \Picardfunctor_{X/B} denotes the Picard functor of X over B
\def\Pic{\mathop{\mathrm{Pic}}\nolimits}
\def\Picardstack{\mathcal{P}\!\mathit{ic}}
\def\Picardfunctor{\mathrm{Pic}}
\def\Deformationcategory{\mathcal{D}\!\mathit{ef}}
\usepackage{cleveref}
\crefname{originallemma}{original lemma}{original lemmas}
\Crefname{originallemma}{Original lemma}{Original lemmas}
\makeatletter
\newlabel{algebra-lemma-lci-CM}{{00SB}{source}{}{sourcepointer-1}{}}
\newlabel{algebra-lemma-lci-CM@cref}{{[originallemma][0][]00SB}{[1][0][]source}{}{}{}}
\makeatother
\crefname{originallemma}{original lemma}{original lemmas}
\Crefname{originallemma}{Original lemma}{Original lemmas}
\makeatletter
\newlabel{divisors-lemma-d-fitting-ideal-omega-smooth}{{0C3K}{source}{}{sourcepointer-2}{}}
\newlabel{divisors-lemma-d-fitting-ideal-omega-smooth@cref}{{[originallemma][0][]0C3K}{[1][0][]source}{}{}{}}
\makeatother
\crefname{originaldefinition}{original definition}{original definitions}
\Crefname{originaldefinition}{Original definition}{Original definitions}
\makeatletter
\newlabel{dpa-definition-lci}{{09Q3}{source}{}{sourcepointer-3}{}}
\newlabel{dpa-definition-lci@cref}{{[originaldefinition][0][]09Q3}{[1][0][]source}{}{}{}}
\makeatother
\crefname{originallemma}{original lemma}{original lemmas}
\Crefname{originallemma}{Original lemma}{Original lemmas}
\makeatletter
\newlabel{dpa-lemma-check-lci-agrees}{{09Q6}{source}{}{sourcepointer-4}{}}
\newlabel{dpa-lemma-check-lci-agrees@cref}{{[originallemma][0][]09Q6}{[1][0][]source}{}{}{}}
\makeatother
\crefname{originalreference}{original reference}{original references}
\Crefname{originalreference}{Original reference}{Original references}
\makeatletter
\newlabel{lemma-nodal}{{0C4D}{source}{}{sourcepointer-5}{}}
\newlabel{lemma-nodal@cref}{{[originalreference][0][]0C4D}{[1][0][]source}{}{}{}}
\makeatother
\crefname{originalreference}{original reference}{original references}
\Crefname{originalreference}{Original reference}{Original references}
\makeatletter
\newlabel{lemma-nodal-algebraic}{{0C49}{source}{}{sourcepointer-6}{}}
\newlabel{lemma-nodal-algebraic@cref}{{[originalreference][0][]0C49}{[1][0][]source}{}{}{}}
\makeatother
\crefname{originaldefinition}{original definition}{original definitions}
\Crefname{originaldefinition}{Original definition}{Original definitions}
\makeatletter
\newlabel{morphisms-definition-syntomic}{{01UC}{source}{}{sourcepointer-7}{}}
\newlabel{morphisms-definition-syntomic@cref}{{[originaldefinition][0][]01UC}{[1][0][]source}{}{}{}}
\makeatother
\crefname{originallemma}{original lemma}{original lemmas}
\Crefname{originallemma}{Original lemma}{Original lemmas}
\makeatletter
\newlabel{morphisms-lemma-algebraic-residue-field-extension-closed-point-fibre}{{01TE}{source}{}{sourcepointer-8}{}}
\newlabel{morphisms-lemma-algebraic-residue-field-extension-closed-point-fibre@cref}{{[originallemma][0][]01TE}{[1][0][]source}{}{}{}}
\makeatother
\crefname{originallemma}{original lemma}{original lemmas}
\Crefname{originallemma}{Original lemma}{Original lemmas}
\makeatletter
\newlabel{morphisms-lemma-local-complete-intersection}{{01UG}{source}{}{sourcepointer-9}{}}
\newlabel{morphisms-lemma-local-complete-intersection@cref}{{[originallemma][0][]01UG}{[1][0][]source}{}{}{}}
\makeatother
\crefname{originallemma}{original lemma}{original lemmas}
\Crefname{originallemma}{Original lemma}{Original lemmas}
\makeatletter
\newlabel{morphisms-lemma-residue-field-quasi-finite}{{01TG}{source}{}{sourcepointer-10}{}}
\newlabel{morphisms-lemma-residue-field-quasi-finite@cref}{{[originallemma][0][]01TG}{[1][0][]source}{}{}{}}
\makeatother
\crefname{originallemma}{original lemma}{original lemmas}
\Crefname{originallemma}{Original lemma}{Original lemmas}
\makeatletter
\newlabel{morphisms-lemma-smooth-syntomic}{{01VD}{source}{}{sourcepointer-11}{}}
\newlabel{morphisms-lemma-smooth-syntomic@cref}{{[originallemma][0][]01VD}{[1][0][]source}{}{}{}}
\makeatother
\crefname{originallemma}{original lemma}{original lemmas}
\Crefname{originallemma}{Original lemma}{Original lemmas}
\makeatletter
\newlabel{morphisms-lemma-unramified-over-field}{{02G7}{source}{}{sourcepointer-12}{}}
\newlabel{morphisms-lemma-unramified-over-field@cref}{{[originallemma][0][]02G7}{[1][0][]source}{}{}{}}
\makeatother
\crefname{originallemma}{original lemma}{original lemmas}
\Crefname{originallemma}{Original lemma}{Original lemmas}
\makeatletter
\newlabel{morphisms-lemma-unramified-quasi-finite}{{02V5}{source}{}{sourcepointer-13}{}}
\newlabel{morphisms-lemma-unramified-quasi-finite@cref}{{[originallemma][0][]02V5}{[1][0][]source}{}{}{}}
\makeatother
\crefname{originallemma}{original lemma}{original lemmas}
\Crefname{originallemma}{Original lemma}{Original lemmas}
\makeatletter
\newlabel{varieties-lemma-geometrically-reduced-dense-smooth-open}{{056V}{source}{}{sourcepointer-14}{}}
\newlabel{varieties-lemma-geometrically-reduced-dense-smooth-open@cref}{{[originallemma][0][]056V}{[1][0][]source}{}{}{}}
\makeatother
\crefname{originallemma}{original lemma}{original lemmas}
\Crefname{originallemma}{Original lemma}{Original lemmas}
\makeatletter
\newlabel{varieties-lemma-geometrically-reduced-upstairs}{{0384}{source}{}{sourcepointer-15}{}}
\newlabel{varieties-lemma-geometrically-reduced-upstairs@cref}{{[originallemma][0][]0384}{[1][0][]source}{}{}{}}
\makeatother
\crefname{originalsection}{original section}{original sections}
\Crefname{originalsection}{Original section}{Original sections}
\makeatletter
\newlabel{varieties-section-algebraic-schemes}{{06LF}{source}{}{sourcepointer-16}{}}
\newlabel{varieties-section-algebraic-schemes@cref}{{[originalsection][0][]06LF}{[1][0][]source}{}{}{}}
\makeatother
\title{At-worst-nodal curves are exactly local complete intersections whose Fitting-ideal subscheme is unramified over the base field}
\author{The Stacks Project Authors}
\date{stacks-project:curves.tex:local-snapshot-20261012; GNU FDL 1.2 or later}
\begin{document}
\maketitle
\noindent\textit{Editorial note.} All delivered mathematical text is the author's own, copied from the pinned Stacks Project chapter file curves.tex; no statement or proof step was rewritten, reordered or removed.\par\medskip
\noindent\textit{Editorial note.} The counted proof is the author's complete argument: both directions with their hinge steps are retained verbatim, including the author's editorial remark that the proof just puts together earlier results and may hide what is really going on.\par\medskip
\noindent\textit{Editorial note.} The label lemma-nodal-lci names the counted statement, and the proof environment immediately following it in the source is delivered verbatim as its proof.\par\medskip
\noindent\textit{Editorial note.} The counted statement is carried by the author's own lemma environment in the chapter source.\par\medskip
\noindent\textit{Allowed standard prerequisites (editorial).}\begin{itemize}
\item Behaviour of the first Fitting ideal subscheme. Divisors, Lemma divisors-lemma-d-fitting-ideal-omega-smooth; Morphisms, Lemmas morphisms-lemma-unramified-quasi-finite, morphisms-lemma-residue-field-quasi-finite, morphisms-lemma-algebraic-residue-field-extension-closed-point-fibre. Identifies the smooth locus with the complement of the Fitting subscheme Z, and the quasi-finiteness of Z -> Spec(k) giving closed points with finite separable residue fields.
\item Local intersection notion comparisons. Algebra, Lemma algebra-lemma-lci-CM; Divided Power Algebra, Definition dpa-definition-lci with Lemma dpa-lemma-check-lci-agrees; Morphisms, Lemmas morphisms-lemma-smooth-syntomic, morphisms-lemma-local-complete-intersection, morphisms-lemma-unramified-over-field. Transfers the complete-intersection property of the completed local rings to the scheme-level lci notion and gives unramifiedness of Z over k.
\end{itemize}\medskip
\ifdefined\hypertarget\hypertarget{sourcepointer-1}{}\fi
\noindent\textit{Source reference.} Original lemma 00SB: sources/stacks-chow/tags.txt, source lines 1033--1033. Prerequisite pointer invoked at proof line(s) 4727; it names a standard fact of a sibling chapter that the counted argument uses, and every step of the counted proof built on it is retained verbatim in the delivered body.\par\medskip
\ifdefined\hypertarget\hypertarget{sourcepointer-2}{}\fi
\noindent\textit{Source reference.} Original lemma 0C3K: sources/stacks-chow/tags.txt, source lines 15195--15195. Prerequisite pointer invoked at proof line(s) 4733; it names a standard fact of a sibling chapter that the counted argument uses, and every step of the counted proof built on it is retained verbatim in the delivered body.\par\medskip
\ifdefined\hypertarget\hypertarget{sourcepointer-3}{}\fi
\noindent\textit{Source reference.} Original definition 09Q3: sources/stacks-chow/tags.txt, source lines 12271--12271. Prerequisite pointer invoked at proof line(s) 4755; it names a standard fact of a sibling chapter that the counted argument uses, and every step of the counted proof built on it is retained verbatim in the delivered body.\par\medskip
\ifdefined\hypertarget\hypertarget{sourcepointer-4}{}\fi
\noindent\textit{Source reference.} Original lemma 09Q6: sources/stacks-chow/tags.txt, source lines 12274--12274. Prerequisite pointer invoked at proof line(s) 4756; it names a standard fact of a sibling chapter that the counted argument uses, and every step of the counted proof built on it is retained verbatim in the delivered body.\par\medskip
\ifdefined\hypertarget\hypertarget{sourcepointer-5}{}\fi
\noindent\textit{Source reference.} Original reference 0C4D: sources/stacks-curves/pinned-original-source.tex, source lines 4391--4391. Prerequisite pointer invoked at proof line(s) 4731; it names a standard fact of the same chapter that the counted argument uses, and every step of the counted proof built on it is retained verbatim in the delivered body.\par\medskip
\ifdefined\hypertarget\hypertarget{sourcepointer-6}{}\fi
\noindent\textit{Source reference.} Original reference 0C49: sources/stacks-curves/pinned-original-source.tex, source lines 4090--4090. Prerequisite pointer invoked at proof line(s) 4753; it names a standard fact of the same chapter that the counted argument uses, and every step of the counted proof built on it is retained verbatim in the delivered body.\par\medskip
\ifdefined\hypertarget\hypertarget{sourcepointer-7}{}\fi
\noindent\textit{Source reference.} Original definition 01UC: sources/stacks-chow/tags.txt, source lines 2405--2405. Prerequisite pointer invoked at proof line(s) 4747; it names a standard fact of a sibling chapter that the counted argument uses, and every step of the counted proof built on it is retained verbatim in the delivered body.\par\medskip
\ifdefined\hypertarget\hypertarget{sourcepointer-8}{}\fi
\noindent\textit{Source reference.} Original lemma 01TE: sources/stacks-chow/tags.txt, source lines 2371--2371. Prerequisite pointer invoked at proof line(s) 4725; it names a standard fact of a sibling chapter that the counted argument uses, and every step of the counted proof built on it is retained verbatim in the delivered body.\par\medskip
\ifdefined\hypertarget\hypertarget{sourcepointer-9}{}\fi
\noindent\textit{Source reference.} Original lemma 01UG: sources/stacks-chow/tags.txt, source lines 2409--2409. Prerequisite pointer invoked at proof line(s) 4757; it names a standard fact of a sibling chapter that the counted argument uses, and every step of the counted proof built on it is retained verbatim in the delivered body.\par\medskip
\ifdefined\hypertarget\hypertarget{sourcepointer-10}{}\fi
\noindent\textit{Source reference.} Original lemma 01TG: sources/stacks-chow/tags.txt, source lines 2373--2373. Prerequisite pointer invoked at proof line(s) 4722; it names a standard fact of a sibling chapter that the counted argument uses, and every step of the counted proof built on it is retained verbatim in the delivered body.\par\medskip
\ifdefined\hypertarget\hypertarget{sourcepointer-11}{}\fi
\noindent\textit{Source reference.} Original lemma 01VD: sources/stacks-chow/tags.txt, source lines 2442--2442. Prerequisite pointer invoked at proof line(s) 4745; it names a standard fact of a sibling chapter that the counted argument uses, and every step of the counted proof built on it is retained verbatim in the delivered body.\par\medskip
\ifdefined\hypertarget\hypertarget{sourcepointer-12}{}\fi
\noindent\textit{Source reference.} Original lemma 02G7: sources/stacks-chow/tags.txt, source lines 3218--3218. Prerequisite pointer invoked at proof line(s) 4751; it names a standard fact of a sibling chapter that the counted argument uses, and every step of the counted proof built on it is retained verbatim in the delivered body.\par\medskip
\ifdefined\hypertarget\hypertarget{sourcepointer-13}{}\fi
\noindent\textit{Source reference.} Original lemma 02V5: sources/stacks-chow/tags.txt, source lines 3756--3756. Prerequisite pointer invoked at proof line(s) 4719; it names a standard fact of a sibling chapter that the counted argument uses, and every step of the counted proof built on it is retained verbatim in the delivered body.\par\medskip
\ifdefined\hypertarget\hypertarget{sourcepointer-14}{}\fi
\noindent\textit{Source reference.} Original lemma 056V: sources/stacks-chow/tags.txt, source lines 6723--6723. Prerequisite pointer invoked at proof line(s) 4740; it names a standard fact of a sibling chapter that the counted argument uses, and every step of the counted proof built on it is retained verbatim in the delivered body.\par\medskip
\ifdefined\hypertarget\hypertarget{sourcepointer-15}{}\fi
\noindent\textit{Source reference.} Original lemma 0384: sources/stacks-chow/tags.txt, source lines 4223--4223. Prerequisite pointer invoked at proof line(s) 4738; it names a standard fact of a sibling chapter that the counted argument uses, and every step of the counted proof built on it is retained verbatim in the delivered body.\par\medskip
\ifdefined\hypertarget\hypertarget{sourcepointer-16}{}\fi
\noindent\textit{Source reference.} Original section 06LF: sources/stacks-chow/tags.txt, source lines 8464--8464. Prerequisite pointer invoked at proof line(s) 4729; it names a standard fact of a sibling chapter that the counted argument uses, and every step of the counted proof built on it is retained verbatim in the delivered body.\par\medskip
\begin{lemma}
\label{lemma-nodal-lci}
Let $k$ be a field. Let $X$ be a locally algebraic $k$-scheme
equidimensional of dimension $1$.
The following are equivalent
\begin{enumerate}
\item the singularities of $X$ are at-worst-nodal, and
\item $X$ is a local complete intersection over $k$
and the closed subscheme $Z \subset X$ cut out by the
first fitting ideal of $\Omega_{X/k}$ is unramified over $k$.
\end{enumerate}
\end{lemma}

\begin{proof}
We urge the reader to find their own proof of
this lemma; what follows is just putting together earlier results
and may hide what is really going on.

\medskip\noindent
Assume (2). Since $Z \to \Spec(k)$ is quasi-finite
(Morphisms, Lemma \ref{morphisms-lemma-unramified-quasi-finite})
we see that the residue fields of points $x \in Z$ are finite
over $k$ (as well as separable) by
Morphisms, Lemma \ref{morphisms-lemma-residue-field-quasi-finite}.
Hence each $x \in Z$ is a closed point of $X$ by
Morphisms, Lemma
\ref{morphisms-lemma-algebraic-residue-field-extension-closed-point-fibre}.
The local ring $\mathcal{O}_{X, x}$ is Cohen-Macaulay by
Algebra, Lemma \ref{algebra-lemma-lci-CM}.
Since $\dim(\mathcal{O}_{X, x}) = 1$ by dimension theory
(Varieties, Section \ref{varieties-section-algebraic-schemes}), we conclude
that $\text{depth}(\mathcal{O}_{X, x})) = 1$. Thus $x$ is a node
by Lemma \ref{lemma-nodal}. If $x \in X$, $x \not \in Z$, then
$X \to \Spec(k)$ is smooth at $x$ by
Divisors, Lemma \ref{divisors-lemma-d-fitting-ideal-omega-smooth}.

\medskip\noindent
Assume (1). Under this assumption $X$ is geometrically reduced
at every closed point (see
Varieties, Lemma \ref{varieties-lemma-geometrically-reduced-upstairs}).
Hence $X \to \Spec(k)$ is smooth on a dense open by
Varieties, Lemma \ref{varieties-lemma-geometrically-reduced-dense-smooth-open}.
Thus $Z$ is closed and consists of closed points.
By Divisors, Lemma \ref{divisors-lemma-d-fitting-ideal-omega-smooth}
the morphism $X \setminus Z \to \Spec(k)$ is smooth.
Hence $X \setminus Z$ is a local complete intersection by
Morphisms, Lemma \ref{morphisms-lemma-smooth-syntomic}
and the definition of a local complete intersection in
Morphisms, Definition \ref{morphisms-definition-syntomic}.
By Lemma \ref{lemma-nodal} for every point $x \in Z$
the local ring $\mathcal{O}_{Z, x}$ is equal to $\kappa(x)$
and $\kappa(x)$ is separable over $k$. Thus $Z \to \Spec(k)$
is unramified (Morphisms, Lemma \ref{morphisms-lemma-unramified-over-field}).
Finally, Lemma \ref{lemma-nodal} via part (3) of
Lemma \ref{lemma-nodal-algebraic}, shows that $\mathcal{O}_{X, x}$
is a complete intersection in the sense of
Divided Power Algebra, Definition \ref{dpa-definition-lci}.
However, Divided Power Algebra, Lemma \ref{dpa-lemma-check-lci-agrees}
and Morphisms, Lemma \ref{morphisms-lemma-local-complete-intersection}
show that this agrees with the notion used to define a local
complete intersection scheme over a field and the proof is complete.
\end{proof}

\end{document}
